\documentclass{beamer}
\usepackage[utf8]{inputenc}
\usepackage[ngerman]{babel}
\usepackage{amsmath,amssymb}
\title{Turing-Maschinen I}
\begin{document}
\begin{frame}{Turing-Maschinen}
\begin{block}{Beispiel: Multiplikation}
Zu zwei natürlichen Zahlen $i_1, i_2 \in \mathbf{N}$ (Eingabe) suchen wir das entsprechende Produkt $i_1 \cdot i_2$ (Ausgabe).

Um die Zahlen $i_1$ und $i_2$ in der Eingabe voneinander trennen zu können, erweitern wir das Alphabet um ein Trennsymbol: $\Sigma = \{0, 1, \#\}$.

Die entsprechende Funktion $f : \Sigma^* \to \Sigma^*$ ist dann gegeben durch
\[ f(\mathrm{bin}(i_1)\#\mathrm{bin}(i_2)) = \mathrm{bin}(i_1 \cdot i_2) \]
\end{block}
Also zum Beispiel $f(110\#10) = 1100$. Was ist $f(\#\#)$?
\end{frame}
\begin{frame}{Mehr Mathe}
Die Laufzeit ist $T(n) \le c \cdot n^2 + \log_2 n$ für alle $n \ge n_0$, und es gilt $x^{k+1} = \sum_{i=0}^{k} a_i$.
Größe, Übergänge, Äpfel, Öfen, Straße.
\[ \delta : Q \times \Gamma \to Q \times \Gamma \times \{L, R, N\} \]
\end{frame}
\end{document}
